\section{Discussion}\label{sec:discussion}

\paragraph{What PARA's step does, and what a step should do.}
PARA's design rests on two good ideas: keep the frozen prompt prediction as a support-invariant anchor, and move away from it by a step that adapts to the support set.
Our analysis supports the first and not the second, at least in BiomedCLIP's frozen space.
The geometric signals in Eq.~\eqref{eq:eta} mostly measure properties of the pretrained model, the gap between modalities and the narrow cone of image embeddings,
so the step becomes a fixed function of $K$ (Section~\ref{sec:finding1}).
The standardized blend then turns that step into a switch on binary tasks (Section~\ref{sec:finding2}).
Neither observation is a flaw in the idea of anchoring; both concern how the step is computed and applied.
What the step should track is the reliability of the trained branch relative to the anchor (Section~\ref{sec:finding3}),
and the blend should let each branch's confidence count (Section~\ref{sec:finding4}).
Neither is solved here.
The anchor's side of the comparison is easy to estimate on the support set; the trained branch's side is not, because the branch is fitted to those same images.
The log-probability blend helps on multi-class tasks but costs worst-case accuracy where the anchor is good, which suggests that the branches' confidences need calibrating before they are blended.

\paragraph{Recommendations for evaluating adaptive blending.}
Four practices would have exposed these effects early, and we suggest them for future work on prompt-anchored or residual adaptation.
\begin{enumerate}
  \item \emph{Check that an adaptive signal adapts.} Report how much the step varies across support draws and whether it tracks the best step per draw.
  A step with a standard deviation near $0.01$ across draws is a schedule, whatever its formula.
  \item \emph{Compare against the simplest baselines on the same trained models:} zero-shot, visual only, and a step that depends only on $K$.
  Here the shot-only schedule beat the adaptive step, and on binary tasks both reduced to one of the two branches.
  \item \emph{Include binary tasks and balanced metrics.} The switch of Section~\ref{sec:finding2} is invisible on multi-class benchmarks,
  and plain accuracy can reward majority-class predictions (DermaMNIST in our setup).
  \item \emph{Report worst-case alongside mean performance, and test rules on data not used to design them.}
  The anchor-aware rule looked convincing on the design datasets and failed on held-out ones;
  its worst case moved further than its mean.
\end{enumerate}

\paragraph{Scope.}
Our claims concern PARA's step and blend as specified in the paper, reimplemented without official code, and studied in few-shot classification.
PARA also reports base-to-novel generalization \citep{kang2026para}, which we do not examine.
Finding~1 follows from two properties of the frozen space, the modality gap and the narrow image cone, which are known for CLIP-style models generally \citep{liang2022mind};
we expect it to carry over, but we tested only BiomedCLIP.
Finding~2 follows from the standardization alone and holds for any backbone with two classes.

\paragraph{Limitations.}
\begin{itemize}
  \item \emph{One training seed and three shot levels.} Each support draw is trained once, so draw-to-draw variation includes optimization noise.
  We use $K \in \{1, 4, 16\}$ only.
  \item \emph{Datasets and resolution.} We use MedMNIST v2 at $128$ or $224$ pixels with MedMNIST splits and fixed test subsets,
  so absolute numbers are not comparable to the paper's benchmark.
  Three held-out datasets is a small test: two are binary, only OrganAMNIST is multi-class, and BreastMNIST has 156 test images.
  \item \emph{Dependence between settings.} The two prompt sets of a dataset share their support draws, so settings are not independent and pooled counts overstate the evidence;
  the per-setting tests are the safer reading.
  With 10 draws per setting, the smallest attainable Wilcoxon $p$ is about $0.002$, which limits how many settings can survive Holm's correction.
  \item \emph{The anchor-aware constant.} $c = 8$ was chosen in-sample on the design datasets.
  We chose it before training on any held-out dataset, but the choice itself was not recorded with a date (Section~\ref{sec:setup}),
  and leave-one-dataset-out by the same CVaR20 criterion picks a different $c$ for each left-out dataset.
  \item \emph{The log-probability blend was motivated by held-out results.} It has no fitted parameters, but the idea came after we saw the binary failure,
  so its binary results are not a clean held-out test.
  \item \emph{Hindsight quantities.} The oracle step, the best points on the $\eta$ curves, and the zero-shot minus visual-only gap all use test labels.
  We use them only to describe what a good step would do, never as rules.
  \item \emph{Baselines are adapted.} The HOSO-style step chooses PARA's $\eta$ on one held-out shot per class; it follows the idea of \citet{vorster2026hoso}
  but is not their method, which learns the ratio by gradient descent.
\end{itemize}

\paragraph{Future work.}
The most useful next step is an estimate of the trained branch's reliability relative to the anchor that needs no test labels,
for example from held-out shots when $K$ allows, or from agreement between the branches on unlabelled images.
A second is calibrated blending: rescaling each branch's confidence on the support set before combining log-probabilities,
which may keep the multi-class gains of Section~\ref{sec:finding4} without the loss in worst case.
Both should be tested on more held-out datasets, more seeds and a second backbone.

\section{Conclusion}\label{sec:conclusion}
PARA keeps a frozen prompt prediction as an anchor and moves away from it by a step computed from the support set.
In BiomedCLIP, that step is set by the number of shots and not by the support set, because its geometric signals reflect the pretrained space rather than the classes;
on binary tasks the standardized blend turns it into a choice between the anchor and the trained branch.
A schedule in $K$ alone does as well or better.
What the step should depend on is how reliable the anchor is relative to the trained branch; a rule based on the anchor alone did not survive held-out testing,
and blending log-probabilities helps on multi-class tasks at a cost in worst-case accuracy.
Measuring the trained branch's reliability, and calibrating both branches before blending, are the open problems for anchored few-shot adaptation.
